% SEGMENT OLP-0350-S01
% Part: lambda-calculus
% Chapter: introduction
% Section: computable-lambda

\documentclass[../../../include/open-logic-section]{subfiles}

\begin{document}

\olfileid{lam}{int}{lrp}
\olsection{கணிக்கத்தக்க சார்புகள் \usetoken{S}{lambda definable}}

\begin{thm}
\ollabel{thm:computable-lambda}
ஒவ்வொரு கணிக்கத்தக்க பகுதிச் சார்பும் !!{lambda definable}.
\end{thm}

\begin{proof}
ஒவ்வொரு பகுதியாகக் கணிக்கத்தக்க சார்பு~$f$ ஐயும் ஒரு லாம்டா சொல்~$F$
ஆல் !!{lambda defined} ஆகக் காட்ட வேண்டும். க்ளீனியின் இயல்புவடிவத் தேற்றப்படி,
ஒவ்வொரு முதன்மை மீள்வரையறைச் சார்பும் ஒரு லாம்டா சொல்லால் !!{lambda defined}
ஆக உள்ளது என்று முதலில் காட்டி, பின்னர் !!{lambda definable} சார்புகள் பொருத்தமான
சேர்ப்புகள் மற்றும் வரம்பற்ற தேடலின் கீழ் மூடியவை என்று காட்டினால் போதும்.
ஒவ்வொரு முதன்மை மீள்வரையறைச் சார்பும் ஒரு லாம்டா சொல்லால் !!{lambda defined}
ஆக உள்ளது என்று காட்ட, தொடக்கச் சார்புகள் !!{lambda definable} என்றும்,
!!{lambda definable} பகுதிச் சார்புகள் சேர்ப்பு, முதன்மை மீள்வரையறை மற்றும்
வரம்பற்ற தேடலின் கீழ் மூடியவை என்றும் காட்டினால் போதும்.
\end{proof}

நிரூபணத்தின் மீதியை மேலும் எளிதாக வாசிக்க, மரபான ஒரு குறிமானத்தைப்
பயன்படுத்துவோம். எடுத்துக்காட்டாக, $Mxyz$ என்பதற்குப் பதிலாக $M(x, y, z)$
என்று எழுதுவோம். இது குறிப்புணர்த்துவதாக இருந்தாலும், வகையிடப்படாத லாம்டா
கலனத்தின் சொற்களுடன் இடஎண்கள் இணைக்கப்படவில்லை என்பதை நினைவில் கொள்ள வேண்டும்;
எனவே அதே சொல்~$M$ க்கு $M(x,y)$ என்றும் $M(x,y,z,w)$ என்றும் எழுதுவது
சம அளவில் பொருளுடையது. ஆனால் இக்குறிமானம் $M$ ஐ மூன்று மாறிகளின் சார்பாக
நாம் நடத்துகிறோம் என்பதைக் காட்டி, வரையறைகளின் நோக்கத்தை மேலும் தெளிவாக்குகிறது.
இதேபோல், ``$M$ ஐ $M = \lambd[x][\lambd[y][\lambd[z][\dots]]]$ ஆல்
வரையறுக்க'' என்பதற்குப் பதிலாக ``$M(x,y,z) = \dots$ ஆல் $M$ ஐ வரையறுக்க''
என்று கூறுவோம்.

\end{document}
